\chapter*{Introduction}
\phantomsection
\addcontentsline{toc}{chapter}{Introduction}
\markboth{INTRODUCTION}{INTRODUCTION}

\section*{Who this book is for}

You can call \code{softmax} and you cannot say why it subtracts its own maximum
first. You have trained a model and you have never derived the gradient it was
following. You read a paper, understand every sentence around the equations,
and skip the equations.

That is not a character flaw and it is not unusual. It is what happens when a
field grows faster than its textbooks, and when the mathematics that field
rests on is taught in an order designed for mathematicians. This book is
written for the engineer in that position: professionally competent, formally
under-equipped, and aware of it.

It assumes \textbf{nothing}. Not algebra, not trigonometry, not calculus. Part~I
begins with arithmetic and means it. If you have a degree in a numerate
subject, that is not a problem --- the Quizzes exist precisely so that you can
skip what you already have --- but the floor is genuinely the floor, because a
book that claims to start from zero and then quietly assumes indices has
started from one.

\section*{What it is for}

Fluency, not rigour. When you have finished it you will be able to derive
scaled dot-product attention, say exactly what the $1/\sqrt{d_k}$ is doing and
what goes wrong without it, read a loss function as an expression rather than a
name, and look at a benchmark table and say whether the difference in it is
real.

You will not be able to prove things. That is a deliberate limitation of the
method and it is stated again, in each place it bites, in a box that names what
is being asserted and where the proof lives. A book that pretends to give you
both in one volume gives you neither.

\section*{Why the format is strange}

This book uses K.~A. Stroud's programmed-learning method, from
\emph{Engineering Mathematics} (1970). The material is entirely different; the
machinery is his, and it is used because it is the only widely deployed format
in mathematics teaching that was ever measured against a control group.

The machinery is simple. The book is a sequence of numbered \emph{frames}.
Nearly every one demands a response, and the answer is at the top of the next
frame, so you are asked to commit before you check. Over one program that
happens forty to seventy times. The effect is that retrieval practice is not
something you have to organise for yourself at the end of a chapter --- it is
the only way the pages can be turned.

The design also grades its own support: worked example, then example with gaps,
then \emph{one for you to do}, then unaided problems. And it sets traps on
purpose. Where a topic has a standard misconception, the text will walk you
into it, let you commit to the wrong answer, then say flatly that it is wrong
and show you why. That is uncomfortable and it is the most valuable thing in
the format, because a misconception that has been elicited and corrected is
gone, and one that has merely been avoided is waiting.

\section*{How the book is organised}

Nine parts, forty-six programs.

\textbf{Part~I --- Foundation} (F1--F13) assumes nothing and is triaged by
Quiz. It is bent towards the payoff rather than towards a school syllabus:
logarithms get a whole program because log-space arithmetic is load-bearing
three parts later, and the chain rule gets the longest program in the part
because backpropagation is the chain rule and nothing else.

\textbf{Part~II --- Number, precision and cost} (P1--P3) comes before linear
algebra, which is a departure from every mathematics curriculum and is
deliberate. Everything after it is arithmetic performed by a finite machine on
a budget. A book that waits four hundred pages to admit that the machine cannot
represent $0.1$ has spent four hundred pages teaching a fiction.

\textbf{Part~III --- Linear algebra} (P4--P10) is the largest part, because it
is the subject this audience uses daily and understands least.

\textbf{Part~IV --- Discrete structures and argument} (P11--P13) covers
counting, graphs, and how to read a theorem without believing more than it
says. It sits before the calculus so that \emph{DAG} and \emph{for all
$\varepsilon$} are defined the first time they are needed.

\textbf{Part~V --- Calculus and automatic differentiation} (P14--P17) takes the
chain rule all the way to what \code{loss.backward()} does.

\textbf{Part~VI --- Optimisation} (P18--P21) is absent from Stroud's first
volume entirely.

\textbf{Part~VII --- Probability and statistics} (P22--P27) goes past the
normal distribution, where Stroud stops, to inference and Bayes --- because
this audience's real statistical job is deciding whether a measured difference
is real, and that is the part no engineering mathematics text covers.

\textbf{Part~VIII --- Information theory} (P28--P30) is where the loss function
comes from.

\textbf{Part~IX --- Assembling it} (P31--P33) spends the whole book on one
architecture, one training run and one evaluation, and introduces no new
mathematics.

\section*{Two rules this book holds itself to}

\begin{note}
\textbf{Every number is computed, not remembered.} Every numeric value printed
here is produced by a script in \code{code/} and pulled into the text
mechanically. The book does not contain the digits; it contains a reference to
them. A value the scripts stop producing prints a visible marker and fails the
build. This exists because a wrong digit in a mathematics book is not a broken
example --- it is something the reader will carry.
\end{note}

\begin{note}
\textbf{What has not been measured is labelled as judgement.} Ten experiments
are specified in these programs. Where one has not been run, the claim it would
support says so, and the table that would hold its numbers stays empty. Empty
tables are not an oversight; they are the mechanism.
\end{note}

\section*{Both languages, one source}

This book exists in Polish and in English, built from one repository. The two
editions share their structure, their figures, their exercise numbering and
their answers; only the prose differs. A check on every build proves that the
two have the same programs, the same frame counts and the same numbers, so a
correction applied to one edition cannot silently miss the other.

Where Polish and English mathematical usage genuinely differ --- $\tg$ against
$\tan$, the decimal comma, the angle brackets Polish tradition uses for a
closed interval, $D^2(X)$ against $\Var(X)$ --- the decision is recorded in
Appendix~B rather than left to a translator's memory, and it is executed by the
typesetting rather than applied by hand.
